% !TEX root = ../main.tex

% 实验记录	
\begin{table}
	\renewcommand\arraystretch{1.7}
	\centering
	\begin{tabularx}{\textwidth}{|X|X|X|X|}
		\hline
		Major： & Physics & Grade： & 2022 \\
		\hline
		Name： & 杨舒云 & Student number： & 22344020\\
		\hline
		Room temperature： & \degree C & Experimental location： &  \\
		\hline
		Student Signature：& \textbf{In Attachment} & Score： &\\
		\hline
		Experiment time：& 2024/11/1 & Teacher's Signature：&\\
		\hline
	\end{tabularx}
\end{table}
\section{APL1-4 电子自旋共振实验 \\ Experimental Record}


%---------------------------------------------------------------------
% 实验过程记录
\subsection{Content, Procedures \& Results}

% 操作步骤
\subsubsection{Operations}
\begin{enumerate}
	\item 观察磁共振信号并测定电子郎德因子
	\begin{enumerate}
		\item 将高斯计探头换成 DPPH 样品，调节双 T 调配器和短路活塞直至出现共振信号。
		\item 调节励磁电源使磁感应强度 $B$ 在 3300 G 至 3338 G 范围内，微调直至信号最大。
		\item 旋转频率计，记录每次信号跳动对应的微波频率 $f$。
		\item 使用公式 $g = \frac{h f}{\mu_B B}$ 计算朗德因子 $g$。
		\item 进行多组测量后，计算平均值和误差，并进行分析讨论。
	\end{enumerate}
	
	\item 估测 DPPH 样品横向弛豫时间
	\begin{enumerate}
		\item 将扫场信号接入示波器CH1，磁共振信号接入CH2，示波器切换到X-Y模式，生成李萨如图。
		\item 调节扫场相位至两个共振峰重合后，切回YT模式，记录信号并保存。
		\item 直接从李萨如图拟合出横向弛豫时间 $T_2$，或记录半高宽 $\Delta t$，并通过扫场幅度 $B'$ 计算 $T_2$。
		\item 备注栏记录是否成功保存信号、数据的平均处理情况及其他实验现象。
	\end{enumerate}
	
	\item 测量波导波长及微波波长
	\begin{enumerate}
		\item 记录谐振腔的多个谐振点位置 $L_1$, $L_2$, $L_3$。
		\item 根据相邻谐振点位置计算波导波长 $\lambda_g$，公式为：$\lambda_g = 2 (L_2 - L_1)$。
		\item 通过公式计算微波波长 $\lambda$：$\lambda = \frac{\lambda_g}{\sqrt{1 - \left( \frac{\lambda_g}{\lambda_c} \right)^2}}$，其中 $\lambda_c$ 为截止波长。
	\end{enumerate}
\end{enumerate}	

% 实验结果
\subsubsection{Display}

The results are shown in \cref{tab:tab1}.

\begin{enumerate}
	\item \begin{table}[h!]
		\centering
		\caption{Examples of table}
		\label{tab:tab1}
		\begin{tabular}{|c|c|c|c|c|c|}
			\hline
			组1/序号i & 1 & 2 & 3 & 4 & 5 \\
			$v_{1i}(m/s)$ & 1.26 & 1.08 & 1.00 & 0.75 & 0.38 \\
			$f_{1i}(Hz)$ & 40073 & 40127 & 40105 & 40088 & 40066 \\
			\hline
			组2/序号i & 1 & 2 & 3 & 4 & 5 \\
			$v_{2i}(m/s)$ & 1.21 & 1.06 & 0.99 & 0.52 & 0.57 \\
			$f_{2i}(Hz)$ & 40143 & 40125 & 40084 & 40080 & 40067 \\
			\hline
			组3/序号i & 1 & 2 & 3 & 4 & 5 \\
			$v_{3i}(m/s)$ & 1.15 & 0.98 & 0.78 & 0.59 & 0.36 \\
			$f_{3i}(Hz)$ & 40135 & 40115 & 40092 & 40070 & 40044 \\
			\hline
		\end{tabular}
	\end{table}		
\end{enumerate}


%---------------------------------------------------------------------
% 原始数据
\subsection{Original Data}
The original data in the experimental notebook is shown in %\cref{} (signed).

See the \textbf{Attachment} section for the clean of the experimental bench desktop (%\cref{}).

Other raw data are shown in %\cref{}.


%---------------------------------------------------------------------
% 问题记录
\subsection{Difficulties}
\begin{enumerate}
	\item 
\end{enumerate}
% ---